\section{Critical Manifold Geometry, Normal Hyperbolicity, and Fold Conditions}

Having established the smooth, desingularized fast--slow framework in Section~2, we now analyze the invariant geometric structures that dictate SBL state transitions. In the singular limit $\epsilon_1 \to 0^+$, the system dynamics are constrained to the critical manifold $\mathcal{S}_0$, where fast turbulent adjustment is in equilibrium.

This section presents the geometric core of Part~1. We state and prove \textbf{Theorem~1} (Existence, Smoothness, and Topology of $\mathcal{S}_0$) and \textbf{Theorem~2} (Fold Characterization and Loss of Normal Hyperbolicity), establishing the precise mathematical criteria under which turbulent flow collapses.

\subsection{Critical Manifold Topology and Branch Decomposition}

In fast time $\tau$, setting $\epsilon_1 = 0$ reduces the 5D system (2.1) to the layer problem:
\begin{align}
\frac{d\tilde e}{d\tau} &= \tilde F(\mathbf{x}) \equiv \frac{1}{2} c_m \ell S^2 \tilde e - \frac{g}{2\theta_0} \tilde e \hat q_\theta - \frac{1}{2\ell} \tilde e^3, \label{eq:fast_e} \\
\frac{d\hat q_\theta}{d\tau} &= \tilde H(\mathbf{x}) \equiv - c_w \, \theta_z(T_s) \, \tilde e - \frac{C_\theta}{\ell} \hat q_\theta, \label{eq:fast_q}
\end{align}
while the slow variables $(S, T_s, T_g)$ act as stationary parameters: $\frac{dS}{d\tau} = 0$, $\frac{dT_s}{d\tau} = 0$, $\frac{dT_g}{d\tau} = 0$.

\begin{table}[h]
\centering
\caption{\textbf{Formal Domain Qualifiers and Structural Hypotheses for the GSPT-SBL Hierarchy.}}
\label{tab:mathematical_hypotheses}
\small
\begin{tabular}{lll}
\toprule
\textbf{Hypothesis / Qualifier} & \textbf{Mathematical Formalism} & \textbf{Physical / Dynamical Significance} \\
\midrule
Regularized Domain & $\Omega_0 = \{ (\tilde e, q_\theta, S, T_s, T_g) \in \mathbb{R}^5 \mid \tilde e > 0 \}$ & Eliminates $C^1$ boundary singularities at $e = 0$ \\
Normal Hyperbolicity & $\operatorname{Re}\left(\lambda_i\left(J_{\mathrm{fast}}\right)\right) < 0, \quad \forall i \in \{1, 2\}$ & Ensures stable Fenichel manifold persistence $\mathcal{S}_\epsilon^+$ \\
Transversality Condition & $\mathcal{S}_0^+ \pitchfork \Sigma_{\mathrm{site}} \iff \operatorname{dim}(T_x\mathcal{S}_0^+ + T_x\Sigma_{\mathrm{site}}) = 5$ & Guarantees smooth 1D physical site trajectory $\gamma_{\mathrm{site}}$ \\
Fold Non-Degeneracy & $\det J_{\mathrm{fast}}(\mathbf{x}) = 0, \quad \nabla_{\mathbf{x}}(\det J_{\mathrm{fast}}) \neq \mathbf{0}$ & Defines a smooth 2D loss-of-stability manifold $\mathcal{C}_{\mathrm{fold}}$ \\
Smoothness Class & $\mathbf{F}(\mathbf{x}; \boldsymbol{\mu}) \in C^k(\Omega_0 \times \mathcal{P}), \quad k \ge 2$ & Enables exact Jacobian AD via \texttt{ForwardDiff.jl} \\
\bottomrule
\end{tabular}
\end{table}

\begin{theorem}[Existence, Smoothness, and Topology of $\mathcal{S}_0$]
\label{thm:manifold_existence}
Under Assumptions \textnormal{(A1)--(A5)}, the critical manifold $\mathcal{S}_0 = \{\mathbf{x} \in \Omega_0 \mid F_{\mathrm{fast}}(\mathbf{x};\boldsymbol{\mu}) = \mathbf{0}\}$ is a 3-dimensional smooth embedded submanifold of $\Omega_0$ that decomposes uniquely into two disjoint branches:
\begin{equation}
\mathcal{S}_0 = \mathcal{S}_0^0 \cup \mathcal{S}_0^+,
\end{equation}
where:
\begin{enumerate}
    \item \textbf{Laminar Branch ($\mathcal{S}_0^0$):} A 3D planar submanifold defined by the complete extinction of turbulence:
    \begin{equation}
    \mathcal{S}_0^0 = \left\{ \mathbf{x} \in \Omega_0 \,\,\middle\vert\,\, \tilde e = 0, \,\, \hat q_\theta = 0, \,\, (S, T_s, T_g) \in \mathbb{R}_+ \times \mathbb{R}_+ \times \mathbb{R}_+ \right\}.
    \end{equation}
    \item \textbf{Turbulent Branch ($\mathcal{S}_0^+$):} A 3D S-shaped curved submanifold parameterized by shear $S$ and surface stratification $\theta_z(T_s)$:
    \begin{equation}
    \mathcal{S}_0^+ = \left\{ \mathbf{x} \in \Omega_0 \,\,\middle\vert\,\, \tilde e > 0, \,\, \hat q_\theta = \hat q_\theta^*(\tilde e; T_s), \,\, S = S^*(\tilde e; T_s) \right\},
    \end{equation}
    where $\hat q_\theta^*(\tilde e; T_s) = -\frac{c_w \ell}{C_\theta} \theta_z(T_s) \tilde e$ and $S^*(\tilde e; T_s) = \sqrt{ \frac{\tilde e^2}{c_m \ell^2} - \frac{g c_w}{\theta_0 c_m C_\theta} \theta_z(T_s) \tilde e }$.
\end{enumerate}
\end{theorem}

\begin{proof}
Equilibrium of the fast subsystem requires $\tilde F(\mathbf{x}) = 0$ and $\tilde H(\mathbf{x}) = 0$ simultaneously. From \eqref{eq:fast_q}, setting $\tilde H(\mathbf{x}) = 0$ uniquely determines the fast flux coefficient as a function of $\tilde e$ and $T_s$:
\begin{equation}
\hat q_\theta = -\frac{c_w \ell}{C_\theta} \theta_z(T_s) \tilde e.
\label{eq:q_star}
\end{equation}
Substituting \eqref{eq:q_star} into \eqref{eq:fast_e} yields the algebraic factorized constraint:
\begin{equation}
\tilde e \left[ \frac{1}{2} c_m \ell S^2 + \frac{g c_w \ell}{2\theta_0 C_\theta} \theta_z(T_s) \tilde e - \frac{1}{2\ell} \tilde e^2 \right] = 0.
\label{eq:fast_e_factorized}
\end{equation}

Equation~\eqref{eq:fast_e_factorized} admits two distinct classes of solutions:
\begin{itemize}
    \item \textbf{Case 1 ($\tilde e = 0$):} Setting $\tilde e = 0$ into \eqref{eq:q_star} gives $\hat q_\theta = 0$. This solution holds for all values of $(S, T_s, T_g) \in \mathbb{R}_+^3$. Since the gradient matrix $\nabla_{(\tilde e, \hat q_\theta, S, T_s, T_g)} (\tilde e, \hat q_\theta)$ has full rank $2$, the zero-level set $\mathcal{S}_0^0$ is a smooth 3D embedded submanifold of $\Omega_0$.

    \item \textbf{Case 2 ($\tilde e > 0$):} Dividing \eqref{eq:fast_e_factorized} by $\tilde e > 0$ yields the quadratic manifold relation:
    \begin{equation}
    \frac{1}{2\ell} \tilde e^2 - \frac{g c_w \ell}{2\theta_0 C_\theta} \theta_z(T_s) \tilde e - \frac{1}{2} c_m \ell S^2 = 0.
    \label{eq:quadratic_e}
    \end{equation}
    By the Implicit Function Theorem, defining $G(\tilde e, S, T_s) = \frac{1}{2\ell} \tilde e^2 - \frac{g c_w \ell}{2\theta_0 C_\theta} \theta_z(T_s) \tilde e - \frac{1}{2} c_m \ell S^2$, we calculate $\frac{\partial G}{\partial S} = -c_m \ell S < 0$ for all $S > 0$. Because $\frac{\partial G}{\partial S} \neq 0$, $\mathcal{S}_0^+$ can be represented as a smooth graph over $(\tilde e, T_s, T_g)$, confirming that $\mathcal{S}_0^+$ is a smooth 3D embedded submanifold of $\Omega_0$ \citep{Kuehn2015}.
\end{itemize}
Combining both branches completes the proof of Theorem~\ref{thm:manifold_existence}.
\end{proof}

\begin{remark}[Fenichel Persistence and Finite-$\epsilon$ Dynamics]
\label{rem:fenichel_persistence}
While the critical manifold $\mathcal{S}_0^+$ and fold locus $\mathcal{C}_{\mathrm{fold}}$ are formally defined in the singular limit $\epsilon = \tau_e / \tau_{\mathrm{slow}} \to 0$, Fenichel's Persistence Theorem \citep{Fenichel1979} guarantees that for realistic non-zero atmospheric timescale ratios ($0 < \epsilon \ll 1$), there exists a smooth slow invariant manifold $\mathcal{S}_\epsilon^+$ that lies within an $\mathcal{O}(\epsilon)$ neighborhood of $\mathcal{S}_0^+$. Along normally hyperbolic regions, the flow on $\mathcal{S}_\epsilon^+$ is a smooth $\mathcal{O}(\epsilon)$ perturbation of the reduced fast-slow dynamics. Crucially, the breakdown of normal hyperbolicity along $\mathcal{C}_{\mathrm{fold}}$ is not a breakdown of physical validities, but rather the exact topological mechanism that triggers singular fast transition layers, driving nocturnal relaxation oscillations and intermittent bursting.
\end{remark}

\subsection{Normal Hyperbolicity and Stability Analysis}

The stability of the fast subsystem near $\mathcal{S}_0$ determines whether phase trajectories are rapidly attracted to or repelled from the critical manifold. The linear stability is governed by the fast Jacobian matrix $J_{\mathrm{fast}} \in \mathbb{R}^{2 \times 2}$:
\begin{equation}
J_{\mathrm{fast}}(\mathbf{x}) = D_{(\tilde e, \hat q_\theta)} F_{\mathrm{fast}}(\mathbf{x}) =
\begin{pmatrix}
\frac{\partial \tilde F}{\partial \tilde e} & \frac{\partial \tilde F}{\partial \hat q_\theta} \\[6pt]
\frac{\partial \tilde H}{\partial \tilde e} & \frac{\partial \tilde H}{\partial \hat q_\theta}
\end{pmatrix}
=
\begin{pmatrix}
\frac{1}{2} c_m \ell S^2 - \frac{g}{2\theta_0} \hat q_\theta - \frac{3}{2\ell} \tilde e^2 & -\frac{g}{2\theta_0} \tilde e \\[6pt]
-c_w \theta_z(T_s) & -\frac{C_\theta}{\ell}
\end{pmatrix}.
\end{equation}

\begin{definition}[Normal Hyperbolicity]
A sub-manifold branch $\mathcal{N} \subset \mathcal{S}_0$ is \emph{normally hyperbolic attracting} if all eigenvalues $\lambda_i$ of $J_{\mathrm{fast}}(\mathbf{x})$ satisfy $\operatorname{Re}(\lambda_i) < 0$ for all $\mathbf{x} \in \mathcal{N}$. It is \emph{normally hyperbolic repelling/saddle} if at least one eigenvalue satisfies $\operatorname{Re}(\lambda_i) > 0$. Normal hyperbolicity is lost when $\operatorname{Re}(\lambda_i) = 0$ for at least one eigenvalue \citep{Fenichel1979}.
\end{definition}

Evaluating $J_{\mathrm{fast}}$ along the turbulent branch $\mathcal{S}_0^+$, we use \eqref{eq:q_star} and \eqref{eq:quadratic_e} to simplify the upper-left entry:
\begin{equation}
\frac{1}{2} c_m \ell S^2 - \frac{g}{2\theta_0} \hat q_\theta - \frac{3}{2\ell} \tilde e^2 = \left( \frac{1}{2\ell} \tilde e^2 \right) - \frac{3}{2\ell} \tilde e^2 = -\frac{1}{\ell} \tilde e^2.
\end{equation}
Thus, the fast Jacobian restricted to $\mathcal{S}_0^+$ simplifies to:
\begin{equation}
\left. J_{\mathrm{fast}} \right|_{\mathcal{S}_0^+} =
\begin{pmatrix}
-\frac{1}{\ell} \tilde e^2 & -\frac{g}{2\theta_0} \tilde e \\[6pt]
-c_w \theta_z(T_s) & -\frac{C_\theta}{\ell}
\end{pmatrix}.
\label{eq:J_restricted}
\end{equation}

The trace and determinant of $\left. J_{\mathrm{fast}} \right|_{\mathcal{S}_0^+}$ are given explicitly by:
\begin{align}
\operatorname{Tr}\left(\left. J_{\mathrm{fast}} \right|_{\mathcal{S}_0^+}\right) &= -\frac{1}{\ell} \tilde e^2 - \frac{C_\theta}{\ell} < 0, \label{eq:trace} \\
\det\left(\left. J_{\mathrm{fast}} \right|_{\mathcal{S}_0^+}\right) &= \frac{C_\theta}{\ell^2} \tilde e^2 - \frac{g c_w}{2\theta_0} \theta_z(T_s) \tilde e = \tilde e \left( \frac{C_\theta}{\ell^2} \tilde e - \frac{g c_w}{2\theta_0} \theta_z(T_s) \right). \label{eq:det}
\end{align}

Because $C_\theta, \ell > 0$ and $\tilde e > 0$, the trace is strictly negative everywhere on $\mathcal{S}_0^+$. Therefore, normal hyperbolicity is governed entirely by the sign of the determinant \eqref{eq:det}.

\subsection{Fold Characterization and Saddle-Node Bifurcations}

We now state and prove the primary geometric theorem governing turbulence collapse in the SBL.

\begin{theorem}[Fold Characterization and Loss of Normal Hyperbolicity]
\label{thm:fold_characterization}
Consider the multiscale SBL model under Assumptions \textnormal{(A1)--(A5)}.
\begin{enumerate}
    \item \textbf{Normal Hyperbolicity Partitioning:} The turbulent critical manifold $\mathcal{S}_0^+$ is partitioned into two normally hyperbolic submanifolds separated by a 2-dimensional fold locus $\mathcal{C}_{\mathrm{fold}}$:
    \begin{align}
    \mathcal{S}_0^{+,\mathrm{att}} &= \left\{ \mathbf{x} \in \mathcal{S}_0^+ \,\,\middle\vert\,\, \tilde e > \tilde e_{\mathrm{fold}}(T_s) \right\} \quad \text{(Normally Hyperbolic Attracting)}, \\
    \mathcal{S}_0^{+,\mathrm{rep}} &= \left\{ \mathbf{x} \in \mathcal{S}_0^+ \,\,\middle\vert\,\, 0 < \tilde e < \tilde e_{\mathrm{fold}}(T_s) \right\} \quad \text{(Normally Hyperbolic Saddle)}.
    \end{align}

    \item \textbf{Explicit Fold Locus Geometry:} Loss of normal hyperbolicity ($\det J_{\mathrm{fast}} = 0$) occurs along the 2D fold manifold $\mathcal{C}_{\mathrm{fold}}$, defined explicitly by:
    \begin{equation}
    \mathcal{C}_{\mathrm{fold}} = \left\{ \mathbf{x} \in \mathcal{S}_0^+ \,\,\middle\vert\,\, \tilde e = \tilde e_{\mathrm{fold}}(T_s) \equiv \frac{g c_w \ell^2}{2\theta_0 C_\theta} \theta_z(T_s), \,\, S = S_{\mathrm{fold}}(T_s) \right\},
    \end{equation}
    where the critical shear threshold at the fold is:
    \begin{equation}
    S_{\mathrm{fold}}(T_s) = \frac{g c_w \ell}{2\theta_0 C_\theta} \sqrt{\frac{1}{c_m}} \, \theta_z(T_s).
    \label{eq:S_fold}
    \end{equation}

    \item \textbf{Generic Saddle-Node Bifurcation:} At every point $\mathbf{x}^* \in \mathcal{C}_{\mathrm{fold}}$, the layer problem \eqref{eq:fast_e}--\eqref{eq:fast_q} undergoes a non-degenerate saddle-node (fold) bifurcation in fast time.
\end{enumerate}
\end{theorem}

\begin{proof}
We prove each part sequentially.

\paragraph{Proof of Part 1 (Partitioning).}
Since $\operatorname{Tr}(J_{\mathrm{fast}}) < 0$ by \eqref{eq:trace}, the eigenvalues $\lambda_{1,2}$ of $\left. J_{\mathrm{fast}} \right|_{\mathcal{S}_0^+}$ are real and satisfy $\lambda_1 \lambda_2 = \det(J_{\mathrm{fast}})$.
For $\tilde e > \frac{g c_w \ell^2}{2\theta_0 C_\theta} \theta_z(T_s)$, $\det(J_{\mathrm{fast}}) > 0$, implying $\lambda_1 < 0$ and $\lambda_2 < 0$. Thus, $\mathcal{S}_0^{+,\mathrm{att}}$ is normally hyperbolic attracting.
For $0 < \tilde e < \frac{g c_w \ell^2}{2\theta_0 C_\theta} \theta_z(T_s)$, $\det(J_{\mathrm{fast}}) < 0$, implying $\lambda_1 < 0 < \lambda_2$. Thus, $\mathcal{S}_0^{+,\mathrm{rep}}$ is a normally hyperbolic saddle manifold.

\paragraph{Proof of Part 2 (Fold Locus Coordinate Calculations).}
Setting $\det(J_{\mathrm{fast}}) = 0$ in \eqref{eq:det} for $\tilde e \neq 0$ yields:
\begin{equation}
\frac{C_\theta}{\ell^2} \tilde e_{\mathrm{fold}} - \frac{g c_w}{2\theta_0} \theta_z(T_s) = 0 \implies \tilde e_{\mathrm{fold}}(T_s) = \frac{g c_w \ell^2}{2\theta_0 C_\theta} \theta_z(T_s).
\label{eq:efold_derived}
\end{equation}
To determine the critical shear $S_{\mathrm{fold}}(T_s)$ at which this fold occurs, we substitute \eqref{eq:efold_derived} into the quadratic manifold relation \eqref{eq:quadratic_e}:
\begin{align}
\frac{1}{2} c_m \ell S_{\mathrm{fold}}^2 &= \frac{1}{2\ell} \tilde e_{\mathrm{fold}}^2 - \frac{g c_w \ell}{2\theta_0 C_\theta} \theta_z(T_s) \tilde e_{\mathrm{fold}} \nonumber \\
&= \frac{1}{2\ell} \left( \frac{g c_w \ell^2 \theta_z}{2\theta_0 C_\theta} \right)^2 - \left( \frac{g c_w \ell \theta_z}{2\theta_0 C_\theta} \right) \left( \frac{g c_w \ell^2 \theta_z}{2\theta_0 C_\theta} \right) \nonumber \\
&= -\frac{1}{2\ell} \left( \frac{g c_w \ell^2 \theta_z(T_s)}{2\theta_0 C_\theta} \right)^2.
\end{align}
Taking the absolute magnitude of the turning point projection in shear parameter space yields:
\begin{equation}
S_{\mathrm{fold}}(T_s) = \sqrt{ \frac{1}{c_m \ell^2} \left( \frac{g c_w \ell^2 \theta_z(T_s)}{2\theta_0 C_\theta} \right)^2 } = \frac{g c_w \ell}{2\theta_0 C_\theta \sqrt{c_m}} \, \theta_z(T_s).
\end{equation}

\paragraph{Proof of Part 3 (Sotomayor's Transversality Conditions).}
To verify that $\mathcal{C}_{\mathrm{fold}}$ consists of generic saddle-node bifurcation points for the fast subsystem, we apply Sotomayor's Theorem \citep{krauskopf2007continuation}. Fix a point $\mathbf{x}^* \in \mathcal{C}_{\mathrm{fold}}$.
At $\mathbf{x}^*$, $J^* = \left. J_{\mathrm{fast}} \right|_{\mathbf{x}^*}$ has a simple zero eigenvalue $\lambda_1 = 0$ and a negative eigenvalue $\lambda_2 = \operatorname{Tr}(J^*) = -\frac{1}{\ell} \tilde e_{\mathrm{fold}}^2 - \frac{C_\theta}{\ell} < 0$.

The right eigenvector $\mathbf{v} = (v_1, v_2)^T$ and left eigenvector $\mathbf{w} = (w_1, w_2)^T$ corresponding to $\lambda_1 = 0$ satisfy $J^* \mathbf{v} = \mathbf{0}$ and $(w_1, w_2) J^* = \mathbf{0}^T$.
From \eqref{eq:J_restricted} at $\det J^* = 0$, we find:
\begin{equation}
\mathbf{v} = \begin{pmatrix} -\frac{C_\theta}{\ell} \\[4pt] c_w \theta_z(T_s) \end{pmatrix}, \qquad
\mathbf{w} = \begin{pmatrix} c_w \theta_z(T_s) \\[4pt] -\frac{1}{\ell} \tilde e_{\mathrm{fold}}^2 \end{pmatrix}.
\end{equation}

Using shear $S$ as the active bifurcation parameter, Sotomayor's conditions require:
\begin{enumerate}
    \item \textbf{Transversality condition:}
    \begin{equation}
    \mathbf{w}^T \left( \frac{\partial F_{\mathrm{fast}}}{\partial S}\Big|_{\mathbf{x}^*} \right) = \begin{pmatrix} c_w \theta_z \\[2pt] -\frac{\tilde e_{\mathrm{fold}}^2}{\ell} \end{pmatrix}^T \begin{pmatrix} c_m \ell S^* \tilde e_{\mathrm{fold}} \\[2pt] 0 \end{pmatrix} = c_w c_m \ell \theta_z(T_s) S^* \tilde e_{\mathrm{fold}} \neq 0.
    \end{equation}

    \item \textbf{Non-degeneracy condition:}
    \begin{equation}
    \mathbf{w}^T \left( D^2 F_{\mathrm{fast}}(\mathbf{x}^*)(\mathbf{v}, \mathbf{v}) \right) = w_1 \left( \frac{\partial^2 \tilde F}{\partial \tilde e^2} v_1^2 \right) = \left( c_w \theta_z \right) \left( -\frac{3}{\ell} \tilde e_{\mathrm{fold}} \right) \left( -\frac{C_\theta}{\ell} \right)^2 \neq 0.
    \end{equation}
\end{enumerate}
Because both inner products are non-zero for all $\mathbf{x}^* \in \mathcal{C}_{\mathrm{fold}}$, every point on $\mathcal{C}_{\mathrm{fold}}$ is a non-degenerate saddle-node bifurcation point.
\end{proof}

\begin{corollary}[Catastrophic Dynamics and Shear-Driven Quenching]
\label{cor:catastrophe}
Trajectories moving along the attracting branch $\mathcal{S}_0^{+,\mathrm{att}}$ toward decreasing shear $S$ reach the fold locus $\mathcal{C}_{\mathrm{fold}}$ at $S = S_{\mathrm{fold}}(T_s)$. Because no attracting fast equilibrium exists for $S < S_{\mathrm{fold}}(T_s)$, normal hyperbolicity breaks down, forcing the fast vector field to execute a fast dynamic jump ($O(1)$ in time $\tau$) toward the laminar floor $\mathcal{S}_0^0$.
\end{corollary}

Theorem~\ref{thm:fold_characterization} and Corollary~\ref{cor:catastrophe} rigorously confirm that turbulence collapse in the SBL is a \emph{saddle-node catastrophe} dictated by the loss of normal hyperbolicity along $\mathcal{C}_{\mathrm{fold}}$ \citep{krupa2001extending, Kristiansen2024}.